% ============================================================
%  第二讲  函子、Knörrer 周期性与热带拉格朗日
% ============================================================
\section{第二讲：函子、Knörrer 周期性与热带拉格朗日}

第一讲末的计算给出了一条具体路线：先在纤维方向把 $\ell$ 沿 U 形弧移动，再沿无穷远处的端把结果接回 $Y$ 中的拉格朗日。本讲把这条路线整理成函子语言，然后给出三种把几何"换成更好算的东西"的办法。

\subsection{\texorpdfstring{加 $w_V$ 之后：不卷绕的版本}{加 w\_V 之后：不卷绕的版本}}
\origpage{5}

回顾上一讲：$H$ 是 $V=\left(\C^*\right)^n$ 或环簇中的超曲面，镜像是 $(Y^{n+1},W=w_0+w_V)$，其上是 Floer 理论的 wrapped Fukaya 范畴 $F(Y,W)$。

先复述裤子那个例子，因为后面所有构造都以它为模型。取
\[
\bigl(\C^3,\ w_0=-z_1z_2z_3\bigr),
\]
它是裤子的镜像。$L_0$ 是沿 U 形弧平行移动得到的拉格朗日，$\ell_0=\left(\R_+\right)^2\subseteq\left(\C^*\right)^2=w_0^{-1}(-1)$。三个面板与 U 形弧见图~\ref{fig:panels}、图~\ref{fig:uarc}。由 Abouzaid--A.，
\[
\hom(L_0,L_0)\simeq CW^\bullet(\ell_0,\ell_0)\oplus CW^\bullet(\ell_0,\ell_0)[-1]
=K\bigl[x_1^{\pm1},x_2^{\pm1}\bigr],\qquad \partial=\cdot\,(1+x_1+x_2),
\]
于是 $H^\bullet\hom(L_0,L_0)\cong K[x_1^{\pm1},x_2^{\pm1}]/(1+x_1+x_2)\cong\hom(\OO,\OO)$。对一般超曲面 $H\subset\left(\C^*\right)^n$，同样的计算给出 $L_0$ 与 $\OO_H$ 的相配，只差一个生成性陈述；对裤子这个陈述是现成的。把这些合起来就得到 HMS。

现在加上 $w_V$，仍然取
\[
\bigl(\C^3,\ w_0+w_V=-z_1z_2z_3+T(z_1+z_2+z_3)\bigr).
\]
构造完全相同，区别在于此时不做卷绕：Hamilton 扰动是小的，只让 $\arg(z)$ 在一个小区间内变化。

$w_0$ 的纤维现在仍是 $\left(\C^*\right)^2$，而 $w_V$ 在纤维上的限制为
\[
T\Bigl(z_1+z_2-\frac{w_0}{z_1z_2}\Bigr),
\]
它是 $\PP^2$ 的镜像。为了看清这一点，只需把 $z_3=-w_0/(z_1z_2)$ 代入 $T(z_1+z_2+z_3)$。

从 $\ell_k\subset w_0^{-1}(-1)=\left(\C^*\right)^2$ 出发，要求它关于 $w_V$ 容许。这样的 $\ell_k$ 对应 $\OO_{\PP^2}(k)$。这里 $k$ 的含义是：绕一圈时 $\arg(z_i)$ 扭转 $2\pi k$，具体地 $\arg z_2$ 转过 $2\pi k$ 而 $\arg z_1$ 固定为 $0$。在 $\FF\bigl(\left(\C^*\right)^2,w_V\bigr)$ 中有
\[
CF(\ell_0,\ell_k)\cong H^\bullet\bigl(\OO_{\PP^2}(k)\bigr).
\]
沿 U 形弧平行移动，得到 $Y$ 中的容许拉格朗日 $L_0,L_k$，并且
\begin{equation}
\hom^\bullet(L_0,L_k)\ \cong\ CF^\bullet(\ell_0,\ell_k)\oplus CF^\bullet(\ell_0,\ell_{k-1})[-1],
\label{eq:cone}
\end{equation}
其中微分为 $\partial=\cdot\,(x_0+x_1+x_2)$；一般地，$\partial$ 就是乘以 $H$ 的定义方程。第二个因子的平移来自 $w_0$ 绕原点的单值群，它在 B 侧对应 $-\otimes\OO(1)$；一般情形下对应 $-\otimes\OO(H)$。

式~\eqref{eq:cone} 的两边在 B 侧对应同一条正合序列
\[
H^0\bigl(H=\{x_0+x_1+x_2=0\},\ \OO(k)\bigr)
=H^0\bigl(\PP^2,\OO(k)\bigr)\big/\bigl(x_0+x_1+x_2\bigr),
\]
这正是把 $\PP^2$ 上的截面限制到超曲面 $H$ 上。这一计算属于 Abouzaid 与作者。

\begin{remark}
式~\eqref{eq:cone} 里那个平移 $[-1]$ 是整讲的枢纽。它不是技术性修正，而是单值群的影子：绕 $w_0=0$ 一圈，纤维上的拉格朗日 $\ell_{k-1}$ 被搬成 $\ell_k$ 的"另一半"。忘掉这一项，$\hom^\bullet(L_0,L_k)$ 就不再与 $H^\bullet(\OO(k))$ 同构。
% [Tutor Note: 易错] 平移项 $[-1]$ 来自单值群，不能省略；它是 Knörrer 周期性在链层面的痕迹。来源：原书第 5 页式 $\hom^\bullet(L_0,L_k)\cong CF^\bullet(\ell_0,\ell_k)\oplus CF^\bullet(\ell_0,\ell_{k-1})[-1]$。
\end{remark}

% [Tutor Note: 易错] $U$ 与 $\cap$ 只是一对伴随函子，并不是互逆等价：$U$ 把 $F$ 中的拉格朗日沿 U 形弧平行移动，搬到 $Y$ 中；$\cap$ 则取 $L$ 在无穷远处的"端点"。正因如此才有非平凡的正合三角 $\mu^{-1}\xrightarrow{S}\id\to\cap U$，而不是等号。来源：原书第 6 页函子表与正合三角。
\subsection{\texorpdfstring{函子 $U$、$\cap$、$\mu$ 与 $S$}{函子 U、cap、mu 与 S}}
\origpage{6}

上一小节的构造可以整体地写成函子。$Y$ 上的 Fukaya 范畴与纤维
\[
F=w_0^{-1}(1)
\]
上的 Fukaya 范畴由若干函子相联系：
\[
\mu\subset F(F,w_V)\ \substack{\xrightarrow{\ U\ }\\[-2pt]\xleftarrow{\ \cap\ }}\ F(Y,w_0+w_V).
\]
两个方向分别是：
\begin{itemize}
  \item $U\ell$：把容许的 $\ell\subset F$ 沿 U 形弧平行移动。今天的约定是把 $U\ell$ 记作 $\ell'$，图上写作 $x\mapsto\mu^{-1}x$。
  \item $\cap L$：$L\subset Y$ 在无穷远处的端。严格说，只有当 $w_0|_L$ 落在 sacred rays 附近时，$\cap L$ 才落在 $\Tw F(F,w_V)$ 中。
\end{itemize}
此外还有两个对象：$\mu$ 是 $w_0$ 绕原点的单值群；$S$ 是计数截面的自然变换
\[
S\colon\ \mu^{-1}\longrightarrow\id ,
\]
它属于 Seidel。具体地，$s_\ell^0\in CF^\bullet\bigl(\mu^{-1}(\ell),\ell\bigr)$ 计数 $w_0\colon Y\to\C$ 在圆盘上的全纯截面。

在 $\FF(F,w_V)$ 上有一个函子的正合三角
\[
\mu^{-1}\xto{\ S\ }\id\xto{\ \ }\cap U ,
\]
这一条属于 Abouzaid--Ganatra；在 $\FF(Y,w_0+w_V)$ 上另有对偶的一个，对应另一个伴随。

\begin{theo}[Abouzaid--A.]\label{thm:square}
设 $H\hookrightarrow V$ 是 $\left(\C^*\right)^n$ 或环 Fano 簇中的超曲面。则下面的 HMS 方块交换：
\[
\begin{array}{ccc}
\mu\subset F(F,w_V) & \substack{\xrightarrow{\ U\ }\\[-2pt]\xleftarrow{\ \cap\ }}\ \ \cong & F(Y,w_0+w_V)\\[4pt]
\Big\downarrow{\scriptstyle\cong\ \hms\ \text{（环簇）}} & & \Big\uparrow{\scriptstyle\cong\ \hms\ \text{（超曲面）}}\\[4pt]
-\otimes\OO(H)\subset \db(V) & \substack{\xrightarrow{\ i^*\ }\\[-2pt]\xleftarrow{\ i_*\ }} & \db(H)
\end{array}
\]
两处竖直箭头都是正合等价。
\end{theo}

\begin{remark}[关键点]\label{rem:keypoint}
由显式计算，自然变换 $S\colon i^*\to\id$ 在 HMS 下对应乘以 $H$ 的定义方程，即
\[
f\colon\ -\otimes\OO(-H)\ \longrightarrow\ \id .
\]
这正是把几何信息（$H$ 的方程）转成代数信息（乘法）的那一步。
\end{remark}

用这套语言，第一讲的计算可以重写为
\[
\hom(U\ell,U\ell')\cong\mathrm{Cone}\Bigl(\hom(\ell,\mu^{-1}\ell')\xto{\ S\ }\hom(\ell,\ell')\Bigr)
\cong\hom(\ell,\cap U\ell'),
\]
它由正合三角经伴随性得到；B 侧对应的式子是
\[
\hom_H(i^*E,i^*E')\cong\mathrm{Cone}\Bigl(\hom_V(E,E'\otimes\OO(-H))\xto{\ f\ }\hom_V(E,E')\Bigr).
\]

\subsection{中心纤维与三条路线}
\origpage{7}

除纤维 $F=w_0^{-1}(1)$ 之外，还可以考虑中心纤维 $F_0=w_0^{-1}(0)$。预期它是补集 $V\setminus H$ 的镜像。理由如下：$\db(V\setminus H)$ 由 $\db(V)$ 沿 $f$ 局部化得到，因此自然可以尝试把 $\FF(F_0,w_V)$ 定义为局部化 $\FF(F,w_V)[S^{-1}]$。若如此，上面的 HMS 陈述就给出
\[
\db(V\setminus H)\simeq\FF(F_0,w_V).
\]
这条思路与 Martin Jeffs 的工作有关。

接下来要找各种办法，把 $(Y,W)$ 的几何换成更友好的东西。三条路线：
\begin{enumerate}[label=(\arabic*)]
  \item Knörrer 周期性（M. Jeffs）；
  \item 纤维 $\bigl(F=\left(\C^*\right)^n,w_V\bigr)$ 中的热带拉格朗日（J. Hicks）；
  \item 曲线的三价构形（A.--Efimov--Orlov），这是下一讲的主题。
\end{enumerate}

\paragraph{路线 (1)：Knörrer 周期性.}\label{sec:knonrer}

回忆第一讲注~\ref{rem:linebundle}。在简单例子中 $Y$ 是某个底簇 $Z$ 上的线丛 $K_Z\to Z$，而 $W=u\cdot\sigma$，其中 $\sigma\in H^0\bigl(Z,K_Z^{-1}\bigr)$。由 Orlov 的 Knörrer 周期性，
\[
\dsg(Y,w_0)\simeq\db\Coh\bigl(\sigma^{-1}(0)\bigr),
\]
右端通常奇异，它是 $(n-1)$ 维的镜像。例如裤子 $\{x_1+x_2+1=0\}$ 对应 $\{z_1z_2=0\}\subset\C^2$（$\sigma$ 是 $\C^2$ 的镜像函数）；类似地 $T_{n-1}$ 对应 $\{\prod_i z_i=0\}\subset\C^n$。对第一讲的例~\ref{ex:o3}，对应的则是 $\{z_0z_1z_2=0\}\subset\PP^2$，其紧化是光滑三次曲线
\[
\{z_0z_1z_2+T\textstyle\sum_i z_i^3=0\}.
\]
亏格更高时这类镜像带叠结构且非既约，需要把 $Y$ 取成环叠，本册不予考虑。

Knörrer 周期性还可以"相对化"。把 $D_0=\sigma^{-1}(0)$ 看成 $Z$ 上一族反典范除子铅笔中的奇异纤维，铅笔的纤维记为 $D_t$。定义
\[
W(D_0):=W(D_t)\bigl[s^{-1}\bigr],
\]
即在计数截面的自然变换 $s\colon\id\to\mu$ 处作局部化。于是可以期待如下的 Knörrer 周期性型结果：
\begin{equation}
W\bigl(\sigma^{-1}(0)\bigr)\ \bigl(:=W(D_t)[s^{-1}]\bigr)\ \cong\ F\Bigl(\frac{K_Z}{Z},\ u\cdot\sigma\Bigr),
\label{eq:jeffs}
\end{equation}
右端是线丛 $K_Z\to Z$ 上的 LG 模型。

\begin{theo}[M. Jeffs，进行中]\label{thm:jeffs}
设 $Z$ 是光滑仿射簇，$f\colon Z\to\C$，$D_0=f^{-1}(0)$ 是奇异超曲面，则
\[
W(D_0):=W(D_t)\bigl[s^{-1}\bigr]
\]
全忠实地嵌入 $\FF(Z\times\C,\ u\cdot f)$，其中 $u$ 是 $\C$ 方向的坐标。
\end{theo}

这个定理应当可以大大推广，但现有证明用的是 GPS 的 sectors 形式，因此需要假设正合性。上面那些例子并不是这个定理特别有趣的用例。

\subsection{热带拉格朗日}
\origpage{8}

\paragraph{路线 (2).} 上一小节把 $Y$ 换成 $Z\times\C$；这一小节把 $Y$ 换成纤维 $\left(\C^*\right)^n$ 本身。对 U 形拉格朗日计算 $HF$，可以理解为下面这条 B 侧式子的 HMS 对应物：
\[
\hom_H(i^*E,i^*E')\cong\mathrm{Cone}\Bigl(\hom_V(E,E'\otimes\OO(-H))\xto{\ f\ }\hom_V(E,E')\Bigr).
\]
但也有一条更直接的路：
\[
\hom_H(i^*E,i^*E')\cong\hom_V(E,i_*i^*E'),
\]
即
\[
\hom_{\FF(Y,W)}(U\ell,U\ell')\cong\hom_{\FF(F,w_V)}(\ell,\cap U\ell').
\]
一般地
\[
\cap U\ell'\simeq\mathrm{Cone}\Bigl(\mu^{-1}(\ell')\xto{\ S\ }\ell'\Bigr)\in\Tw\FF(F,w_V),
\]
而常常能找出显式的代表元。对线丛的镜像（即容许拉格朗日截面），这正是 Jeff Hicks 学位论文中热带拉格朗日的构造。

具体地：
\begin{itemize}
  \item $\OO_V$ 的镜像是 $\Log\colon\left(\C^*\right)^n\to\R^n$ 的零截面，$\ell_0=\left(\R_+\right)^n\subset\left(\C^*\right)^n$；
  \item $\OO_V(-H)$ 的镜像是容许拉格朗日截面 $\ell_{(-H)}\subset\left(\C^*\right)^n$，它由光滑化 $d\varphi=d\Trop f$ 的图像得到，位于
  \[
  T^*\R^n\big/ T^*_\Z\ \cong\ \left(\C^*\right)^n
  \]
  中。
\end{itemize}
以裤子为例，$\varphi=\max(0,\xi_1,\xi_2)$，对应的热带曲线见图~\ref{fig:tropcurve}。

\begin{figure}[htbp]
\centering
\begin{tikzpicture}[scale=1.0,>=Stealth]
  \draw[->] (-2.0,0) -- (2.4,0) node[right] {\footnotesize $\xi_1$};
  \draw[->] (0,-1.9) -- (0,2.3) node[above] {\footnotesize $\xi_2$};
  \draw[very thick,red!70!black] (0,0) -- (2.0,0);
  \draw[very thick,red!70!black] (0,0) -- (0,2.0);
  \draw[very thick,red!70!black] (0,0) -- (-2.0,-2.0);
  \node at (1.3,0.25) {\footnotesize $d\xi_1\ (\equiv 0\ \mathrm{mod}\ T^*_\Z)$};
  \node at (0.35,-0.75) {\footnotesize $d\xi_2$};
  \fill (0,0) circle (1.6pt);
  \node[below left] at (0,0) {\footnotesize $0$};
\end{tikzpicture}
\caption{裤子例中 $d\varphi=d\Trop f$ 的图像：三条射线上的斜率分别为 $d\xi_1$、$d\xi_2$ 与 $0$（模 $T^*_\Z$）。（原稿第 8 页）}
\label{fig:tropcurve}
\end{figure}

有了 $\ell_0$ 与 $\ell_{(-H)}$，就可以做拉格朗日手术；更准确地说，取两者的对称差。结果是"热带拉格朗日"
\[
\ell^H\subset\left(\C^*\right)^n ,
\]
它仍然关于 $w_V$ 单项式容许，并且
\[
\Log\bigl(\ell^H\bigr)\subset\ \text{由 }\varphi\text{ 定义的热带超曲面的一个邻域}.
\]
于是例如
\[
\hom_{\FF(Y,W)}(U\ell_0,U\ell_0)\cong\hom_{\FF((\C^*)^n,w_V)}\bigl(\ell_0,\ell^H\bigr),
\]
对其他线丛类似。

\begin{remark}
这里有一个必须小心的地方：$\ell^H$ 给出的是"某个"以 $\varphi$ 为热带化的超曲面的镜像，而要得到 $\OO_H$（即指定了具体超曲面 $H$）的镜像，细节上还需要额外小心。热带化只记录首项，$H$ 的高阶项丢失了。
% [Tutor Note: 易错] 热带拉格朗日只确定超曲面的热带化，不唯一确定 $H$。要得到 $\OO_H$ 本身，必须把"被丢掉的高阶项"补回去；原稿在此处画了警告记号。来源：原书第 8 页末段。
\end{remark}

\subsection{顶针与消没圈}
\origpage{9}

上一小节的结果可以换个角度看。$\hom_{\FF((\C^*)^n,w_V)}(\ell_0,\ell^H)$ 其实是 $\hom(\ell^H,\ell^H)$ 的一半，即其次数为 $0$ 的那部分；与之对应，
\[
H^0(V,\OO_H)\cong\hom^0(\OO_H,\OO_H).
\]
此外，J. Hicks 证明存在一个拉格朗日配边，实现"对称差"手术，因此 $Y$ 中存在容许拉格朗日，其 $\Log$ 投影的形状是：$\ell_{(H)}$ 与 $\ell_0$ 由一段细颈相连，两端各是一个圆。

这件事为什么重要？因为可以这样理解：$L_0=U\ell_0$ 是"顶针"（thimble），而 $\cap U\ell_0\simeq\ell^H$ 是"消没圈"（vanishing cycle）。一般的图景如下。

设 $Y\xrightarrow{\ \pi\ }\C$ 是 Morse--Bott 纤维化，临界轨迹 $Z$ 光滑且余维 2。取 $\ell\subset Z$，沿平行移动可以得到一个"顶针" $T\ell\subset Y$，它是 $\ell$ 上的 $D^2$-丛。这样得到全忠实嵌入
\[
\FF(Z)\hookrightarrow\FF(Y,\pi),
\]
并且
\[
\hom_{\FF(Y,\pi)}(T\ell,T\ell')\cong\hom_{\FF(Z)}(\ell,\ell')
\hookrightarrow\hom_{S(\text{fiber})}(\cap T\ell,\cap T\ell'),
\]
右端是纤维上的消没圈，即 $S^1$-丛除以 $\ell$。

不过，我们更愿意直接在 $\crit(W)\subset Y$ 上工作，即便超出 Morse--Bott 的情形。

\begin{example}\label{ex:one}
$H=\{-1\}\subset V=\C^*$，$\;(Y,W)=\bigl(\C^2,-z_1z_2\bigr)$。此时 $\FF(Y,W)$ 中对应 $\OO_H$ 的容许拉格朗日 $L_0$ 可以是图~\ref{fig:three-L0} 中任意一个：U 形的 $L_0=U\ell_0$、热带拉格朗日、或 Lefschetz 顶针。
\end{example}

\begin{figure}[htbp]
\centering
\begin{tikzpicture}[scale=0.85,>=Stealth]
  % (a) U 形
  \begin{scope}
    \foreach \x in {0,0.5,1.0}{
      \draw[thick] (\x,0) ellipse (0.18 and 0.55);
    }
    \draw[thick] (-0.1,0) -- (1.1,0);
    \draw[blue!70!black,thick] (-0.05,-0.3) -- (0.25,-0.75) .. controls (0.5,-0.95) and (0.6,-0.95) .. (0.85,-0.75) -- (1.15,-0.3);
    \node at (0.5,-1.35) {\footnotesize $L_0=U\ell_0$（U 形）};
  \end{scope}
  \node at (2.1,-0.35) {\footnotesize $\sim$};
  % (b) 热带拉格朗日
  \begin{scope}[xshift=3.3cm]
    \foreach \x in {0,0.5,1.0}{
      \draw[thick] (\x,0) ellipse (0.18 and 0.55);
    }
    \draw[thick] (-0.1,0) .. controls (0.2,-0.35) and (0.35,0.35) .. (0.55,0)
      .. controls (0.75,-0.35) and (0.9,0.35) .. (1.1,0);
    \node at (0.5,-1.35) {\footnotesize 热带拉格朗日};
  \end{scope}
  \node at (5.4,-0.35) {\footnotesize $\sim$};
  % (c) 顶针
  \begin{scope}[xshift=6.6cm]
    \draw[thick] (0,0.6) .. controls (0.5,1.0) and (0.9,0.6) .. (1.0,0.15);
    \draw[thick] (0,0.6) .. controls (0.5,0.2) and (0.9,-0.2) .. (1.0,0.15);
    \node at (0.5,0.15) {\footnotesize $\times$};
    \draw[thick] (1.0,0.15) -- (2.1,0.15);
    \node at (1.0,-0.95) {\footnotesize Lefschetz 顶针};
  \end{scope}
\end{tikzpicture}
\caption{$\bigl(\C^2,-z_1z_2\bigr)$ 中三种等价的容许拉格朗日：U 形、热带拉格朗日、Lefschetz 顶针。（原稿第 9 页）}
\label{fig:three-L0}
\end{figure}

\begin{example}[裤子]\label{ex:pants-thimble}
$H=\{1+x_1+x_2=0\}\subset V=\left(\C^*\right)^2$，$\;(Y,W)=\bigl(\C^3,-z_1z_2z_3\bigr)$。对 $x\in H$，$\OO_x$ 可以实现为顶针 $L_x$（一个实心环），做法来自 Aganagic--Vafa：取 $S^1\subset\crit W=\cup_3\C$，再加上局部系统。这个 $L_x$ 是有阻碍的，它边界全纯圆盘，圆盘位于 $\C$ 中，面积为 $\val(x)$，落在图~\ref{fig:wedge} 所示的楔形区域里。因此需要给它配一个界定余链，由后者修正无穷远处的 holonomy。
\end{example}

\begin{figure}[htbp]
\centering
\begin{tikzpicture}[scale=1.0,>=Stealth]
  \fill[blue!12] (0,0) -- (40:2.2) -- (80:2.2) -- cycle;
  \draw[->] (-0.3,0) -- (2.5,0) node[right] {\footnotesize $\R^2$};
  \draw[->] (0,-0.3) -- (0,2.4);
  \draw[thick,blue!70!black] (0,0) -- (40:2.2);
  \draw[thick,blue!70!black] (0,0) -- (80:2.2);
  \node[blue!70!black] at (60:2.45) {\footnotesize $\mathrm{area}\mapsto\val(x)$};
\end{tikzpicture}
\caption{有阻碍的顶针所边界圆盘的面积取值区域：落在 $\R^2$ 中由 $\val(x)$ 决定的楔形里。（原稿第 9 页）}
\label{fig:wedge}
\end{figure}

界定余链的选取并不唯一，它对应 Ag-Vafa 与 11 维构造中所谓的 framing。

\subsection{有阻碍的情形与 SYZ 纤维化}
\origpage{10}

上一小节的结论可以算得更具体。设膜落在 $z_1$ 轴上。此时第一个因子中的位置与 holonomy 为
\[
T^{*\arg}\mathrm{hol}(0)=x_1,\qquad \val(x_1)>0 ,
\]
第二个因子是 $S^1_{z_2}$，它在 $L_x$ 中边界一个圆盘。界定余链把无穷远处的有效 holonomy 从 $-1$ 改为
\[
-1-x_1=x_2 ,
\]
其中用到了 spin 结构的变号。于是
\[
\cap L_x\cong T_x=(x_1,x_2).
\]
注意 $x_1+x_2=-1$，这正与 $H$ 的方程一致。另外，
\[
L_x\oplus L_x[1]\cong UT_x ,
\]
这与 $i^*i_*\OO_x\cong\OO_x\oplus\OO_x[1]$ 相符。

对摩天层 $\OO_x$，其实并不存在真正的顶针。不过可以观察到一个分片光滑的替代物。考虑
\[
\cup_3\bigl(\R_{\ge0}\subset\C\bigr)\subset\{z_1z_2z_3=0\},
\]
这条三价图沿平行移动可以加粗成 $Y$ 中的一个分片光滑拉格朗日：$\R_+$ 的每个点给出一个 $S^1$，而 $0$ 处的纤维是 $T^2$（图~\ref{fig:thicken}）。

\begin{figure}[htbp]
\centering
\begin{tikzpicture}[scale=0.9,>=Stealth]
  % 三价图
  \begin{scope}
    \draw[thick] (0,0) -- (150:1.1);
    \draw[thick] (0,0) -- (30:1.1);
    \draw[thick] (0,0) -- (270:1.1);
    \fill (0,0) circle (1.6pt);
    \node at (0,-1.6) {\footnotesize $\cup_3\bigl(\R_{\ge0}\subset\C\bigr)\subset\{z_1z_2z_3=0\}$};
  \end{scope}
  \draw[->,thick] (1.6,-0.2) -- (2.6,-0.2);
  % 加粗
  \begin{scope}[xshift=3.4cm]
    \draw[thick] (0,0) -- (150:1.1);
    \draw[thick] (0,0) -- (30:1.1);
    \draw[thick] (0,0) -- (270:1.1);
    \foreach \a in {150,30,270}{
      \draw[thick] (\a:0.35) ellipse (0.12 and 0.30);
    }
    \fill (0,0) circle (1.6pt);
    \node[align=left] at (1.5,1.35) {\footnotesize 每个 $\R_+$ 的点给出 $S^1$};
    \node[align=left] at (1.5,-1.15) {\footnotesize $0$ 处纤维为 $T^2$};
  \end{scope}
\end{tikzpicture}
\caption{把 $\cup_3(\R_{\ge0}\subset\C)$ 沿平行移动加粗：得到 $Y$ 中分片光滑的拉格朗日。（原稿第 10 页）}
\label{fig:thicken}
\end{figure}

用它在 $\C^3$ 中构造分片光滑的拉格朗日并不难；可以设想用 J. Hicks 的热带拉格朗日技术把这个分片光滑的对象光滑化，但目前还没有人做出来。

不过可以先把问题放在一边，想象直接在 $\crit(W)$ 上做辛几何，或者在 $\bigl(\crit(w_0),w_V\bigr)$ 上做。这是下一讲的主题。作为一个引子：$\C^3\setminus\{z_1z_2z_3=1\}$ 上有 SYZ 纤维化，底空间是 $\R^3$，沿一条三价图奇异；于是可以问，能否为某些分支找到热带拉格朗日。
